% Options for packages loaded elsewhere
\PassOptionsToPackage{unicode}{hyperref}
\PassOptionsToPackage{hyphens}{url}
\documentclass[
]{article}
\usepackage{xcolor}
\usepackage[margin=1in]{geometry}
\usepackage{amsmath,amssymb}
\setcounter{secnumdepth}{-\maxdimen} % remove section numbering
\usepackage{iftex}
\ifPDFTeX
  \usepackage[T1]{fontenc}
  \usepackage[utf8]{inputenc}
  \usepackage{textcomp} % provide euro and other symbols
\else % if luatex or xetex
  \usepackage{unicode-math} % this also loads fontspec
  \defaultfontfeatures{Scale=MatchLowercase}
  \defaultfontfeatures[\rmfamily]{Ligatures=TeX,Scale=1}
\fi
\usepackage{lmodern}
\ifPDFTeX\else
  % xetex/luatex font selection
\fi
% Use upquote if available, for straight quotes in verbatim environments
\IfFileExists{upquote.sty}{\usepackage{upquote}}{}
\IfFileExists{microtype.sty}{% use microtype if available
  \usepackage[]{microtype}
  \UseMicrotypeSet[protrusion]{basicmath} % disable protrusion for tt fonts
}{}
\makeatletter
\@ifundefined{KOMAClassName}{% if non-KOMA class
  \IfFileExists{parskip.sty}{%
    \usepackage{parskip}
  }{% else
    \setlength{\parindent}{0pt}
    \setlength{\parskip}{6pt plus 2pt minus 1pt}}
}{% if KOMA class
  \KOMAoptions{parskip=half}}
\makeatother
\usepackage{color}
\usepackage{fancyvrb}
\newcommand{\VerbBar}{|}
\newcommand{\VERB}{\Verb[commandchars=\\\{\}]}
\DefineVerbatimEnvironment{Highlighting}{Verbatim}{commandchars=\\\{\}}
% Add ',fontsize=\small' for more characters per line
\usepackage{framed}
\definecolor{shadecolor}{RGB}{248,248,248}
\newenvironment{Shaded}{\begin{snugshade}}{\end{snugshade}}
\newcommand{\AlertTok}[1]{\textcolor[rgb]{0.94,0.16,0.16}{#1}}
\newcommand{\AnnotationTok}[1]{\textcolor[rgb]{0.56,0.35,0.01}{\textbf{\textit{#1}}}}
\newcommand{\AttributeTok}[1]{\textcolor[rgb]{0.13,0.29,0.53}{#1}}
\newcommand{\BaseNTok}[1]{\textcolor[rgb]{0.00,0.00,0.81}{#1}}
\newcommand{\BuiltInTok}[1]{#1}
\newcommand{\CharTok}[1]{\textcolor[rgb]{0.31,0.60,0.02}{#1}}
\newcommand{\CommentTok}[1]{\textcolor[rgb]{0.56,0.35,0.01}{\textit{#1}}}
\newcommand{\CommentVarTok}[1]{\textcolor[rgb]{0.56,0.35,0.01}{\textbf{\textit{#1}}}}
\newcommand{\ConstantTok}[1]{\textcolor[rgb]{0.56,0.35,0.01}{#1}}
\newcommand{\ControlFlowTok}[1]{\textcolor[rgb]{0.13,0.29,0.53}{\textbf{#1}}}
\newcommand{\DataTypeTok}[1]{\textcolor[rgb]{0.13,0.29,0.53}{#1}}
\newcommand{\DecValTok}[1]{\textcolor[rgb]{0.00,0.00,0.81}{#1}}
\newcommand{\DocumentationTok}[1]{\textcolor[rgb]{0.56,0.35,0.01}{\textbf{\textit{#1}}}}
\newcommand{\ErrorTok}[1]{\textcolor[rgb]{0.64,0.00,0.00}{\textbf{#1}}}
\newcommand{\ExtensionTok}[1]{#1}
\newcommand{\FloatTok}[1]{\textcolor[rgb]{0.00,0.00,0.81}{#1}}
\newcommand{\FunctionTok}[1]{\textcolor[rgb]{0.13,0.29,0.53}{\textbf{#1}}}
\newcommand{\ImportTok}[1]{#1}
\newcommand{\InformationTok}[1]{\textcolor[rgb]{0.56,0.35,0.01}{\textbf{\textit{#1}}}}
\newcommand{\KeywordTok}[1]{\textcolor[rgb]{0.13,0.29,0.53}{\textbf{#1}}}
\newcommand{\NormalTok}[1]{#1}
\newcommand{\OperatorTok}[1]{\textcolor[rgb]{0.81,0.36,0.00}{\textbf{#1}}}
\newcommand{\OtherTok}[1]{\textcolor[rgb]{0.56,0.35,0.01}{#1}}
\newcommand{\PreprocessorTok}[1]{\textcolor[rgb]{0.56,0.35,0.01}{\textit{#1}}}
\newcommand{\RegionMarkerTok}[1]{#1}
\newcommand{\SpecialCharTok}[1]{\textcolor[rgb]{0.81,0.36,0.00}{\textbf{#1}}}
\newcommand{\SpecialStringTok}[1]{\textcolor[rgb]{0.31,0.60,0.02}{#1}}
\newcommand{\StringTok}[1]{\textcolor[rgb]{0.31,0.60,0.02}{#1}}
\newcommand{\VariableTok}[1]{\textcolor[rgb]{0.00,0.00,0.00}{#1}}
\newcommand{\VerbatimStringTok}[1]{\textcolor[rgb]{0.31,0.60,0.02}{#1}}
\newcommand{\WarningTok}[1]{\textcolor[rgb]{0.56,0.35,0.01}{\textbf{\textit{#1}}}}
\usepackage{longtable,booktabs,array}
\newcounter{none} % for unnumbered tables
\usepackage{calc} % for calculating minipage widths
% Correct order of tables after \paragraph or \subparagraph
\usepackage{etoolbox}
\makeatletter
\patchcmd\longtable{\par}{\if@noskipsec\mbox{}\fi\par}{}{}
\makeatother
% Allow footnotes in longtable head/foot
\IfFileExists{footnotehyper.sty}{\usepackage{footnotehyper}}{\usepackage{footnote}}
\makesavenoteenv{longtable}
\usepackage{graphicx}
\makeatletter
\newsavebox\pandoc@box
\newcommand*\pandocbounded[1]{% scales image to fit in text height/width
  \sbox\pandoc@box{#1}%
  \Gscale@div\@tempa{\textheight}{\dimexpr\ht\pandoc@box+\dp\pandoc@box\relax}%
  \Gscale@div\@tempb{\linewidth}{\wd\pandoc@box}%
  \ifdim\@tempb\p@<\@tempa\p@\let\@tempa\@tempb\fi% select the smaller of both
  \ifdim\@tempa\p@<\p@\scalebox{\@tempa}{\usebox\pandoc@box}%
  \else\usebox{\pandoc@box}%
  \fi%
}
% Set default figure placement to htbp
\def\fps@figure{htbp}
\makeatother
\setlength{\emergencystretch}{3em} % prevent overfull lines
\providecommand{\tightlist}{%
  \setlength{\itemsep}{0pt}\setlength{\parskip}{0pt}}
\usepackage{bookmark}
\IfFileExists{xurl.sty}{\usepackage{xurl}}{} % add URL line breaks if available
\urlstyle{same}
\hypersetup{
  pdftitle={test markdown},
  hidelinks,
  pdfcreator={LaTeX via pandoc}}

\title{test markdown}
\author{}
\date{\vspace{-2.5em}2026-05-01}

\begin{document}
\maketitle

\begin{Shaded}
\begin{Highlighting}[]
\FunctionTok{library}\NormalTok{(dplyr)}
\end{Highlighting}
\end{Shaded}

\begin{verbatim}
## 
## Attaching package: 'dplyr'
\end{verbatim}

\begin{verbatim}
## The following objects are masked from 'package:stats':
## 
##     filter, lag
\end{verbatim}

\begin{verbatim}
## The following objects are masked from 'package:base':
## 
##     intersect, setdiff, setequal, union
\end{verbatim}

\begin{Shaded}
\begin{Highlighting}[]
\FunctionTok{library}\NormalTok{(tidyverse)}
\end{Highlighting}
\end{Shaded}

\begin{verbatim}
## -- Attaching core tidyverse packages ------------------------ tidyverse 2.0.0 --
## v forcats   1.0.1     v readr     2.2.0
## v ggplot2   4.0.3     v stringr   1.6.0
## v lubridate 1.9.5     v tibble    3.3.1
## v purrr     1.2.2     v tidyr     1.3.2
\end{verbatim}

\begin{verbatim}
## -- Conflicts ------------------------------------------ tidyverse_conflicts() --
## x dplyr::filter() masks stats::filter()
## x dplyr::lag()    masks stats::lag()
## i Use the conflicted package (<http://conflicted.r-lib.org/>) to force all conflicts to become errors
\end{verbatim}

\begin{Shaded}
\begin{Highlighting}[]
\FunctionTok{library}\NormalTok{(glmnet)}
\end{Highlighting}
\end{Shaded}

\begin{verbatim}
## Loading required package: Matrix
## 
## Attaching package: 'Matrix'
## 
## The following objects are masked from 'package:tidyr':
## 
##     expand, pack, unpack
## 
## Loaded glmnet 4.1-10
\end{verbatim}

\begin{Shaded}
\begin{Highlighting}[]
\FunctionTok{library}\NormalTok{(sf)  }
\end{Highlighting}
\end{Shaded}

\begin{verbatim}
## Linking to GEOS 3.14.1, GDAL 3.8.5, PROJ 9.5.1; sf_use_s2() is TRUE
\end{verbatim}

\begin{Shaded}
\begin{Highlighting}[]
\FunctionTok{library}\NormalTok{(spatstat) }
\end{Highlighting}
\end{Shaded}

\begin{verbatim}
## Loading required package: spatstat.data
## Loading required package: spatstat.univar
## spatstat.univar 3.1-7
## Loading required package: spatstat.geom
## spatstat.geom 3.7-3
## Loading required package: spatstat.random
## spatstat.random 3.4-5
## Loading required package: spatstat.explore
## Loading required package: nlme
## 
## Attaching package: 'nlme'
## 
## The following object is masked from 'package:dplyr':
## 
##     collapse
## 
## spatstat.explore 3.8-0
## Loading required package: spatstat.model
## Loading required package: rpart
## spatstat.model 3.7-0
## Loading required package: spatstat.linnet
## spatstat.linnet 3.5-0
## 
## spatstat 3.6-0 
## For an introduction to spatstat, type 'beginner'
\end{verbatim}

\begin{Shaded}
\begin{Highlighting}[]
\FunctionTok{library}\NormalTok{(ggplot2)  }
\FunctionTok{library}\NormalTok{(leaflet)}
\FunctionTok{library}\NormalTok{(maps)}
\end{Highlighting}
\end{Shaded}

\begin{verbatim}
## 
## Attaching package: 'maps'
## 
## The following object is masked from 'package:purrr':
## 
##     map
\end{verbatim}

\begin{Shaded}
\begin{Highlighting}[]
\FunctionTok{library}\NormalTok{(tidyverse)}
\FunctionTok{library}\NormalTok{(skimr)}
\FunctionTok{library}\NormalTok{(gridExtra)}
\end{Highlighting}
\end{Shaded}

\begin{verbatim}
## 
## Attaching package: 'gridExtra'
## 
## The following object is masked from 'package:dplyr':
## 
##     combine
\end{verbatim}

\begin{Shaded}
\begin{Highlighting}[]
\FunctionTok{library}\NormalTok{(knitr)}
\end{Highlighting}
\end{Shaded}

\section{Read in Dataset and Perform Summary
Statistics}\label{read-in-dataset-and-perform-summary-statistics}

\begin{Shaded}
\begin{Highlighting}[]
\NormalTok{blowdown}\OtherTok{=}\FunctionTok{read.csv}\NormalTok{(}\StringTok{"/Users/sameehahossain/Desktop/R Files/ef\_w\_blowdown\_ons\_ng\_trns\_pl.csv"}\NormalTok{)}

\CommentTok{\#view column names}
\FunctionTok{colnames}\NormalTok{(blowdown)}
\end{Highlighting}
\end{Shaded}

\begin{verbatim}
##  [1] "facility_name"       "facility_id"         "industry_segment"   
##  [4] "number_of_blowdowns" "reporting_year"      "table_desc"         
##  [7] "table_num"           "total_ch4_emissions" "total_co2_emissions"
## [10] "us_state"
\end{verbatim}

\begin{Shaded}
\begin{Highlighting}[]
\CommentTok{\#Data Quality / Structure (any data is missing or needs cleaning?)}
\FunctionTok{summary}\NormalTok{(blowdown)}
\end{Highlighting}
\end{Shaded}

\begin{verbatim}
##    facility_name   facility_id       industry_segment number_of_blowdowns
##  Length   :1388   Min.   :1010455   Length   :1388    Min.   :   0.00    
##  N.unique :  56   1st Qu.:1012291   N.unique :   1    1st Qu.:   4.00    
##  N.blank  :   0   Median :1012473   N.blank  :   0    Median :  12.00    
##  Min.nchar:   9   Mean   :1012553   Min.nchar:  57    Mean   :  60.31    
##  Max.nchar:  75   3rd Qu.:1012701   Max.nchar:  57    3rd Qu.:  38.00    
##                   Max.   :1014613                     Max.   :3386.00    
##  reporting_year     table_desc       table_num    total_ch4_emissions
##  Min.   :2016   Length   :1388   Length   :1388   Min.   :    0.00   
##  1st Qu.:2018   N.unique :   1   N.unique :   1   1st Qu.:   11.93   
##  Median :2020   N.blank  :   0   N.blank  :   0   Median :  164.00   
##  Mean   :2020   Min.nchar:  58   Min.nchar:  17   Mean   :  620.98   
##  3rd Qu.:2022   Max.nchar:  58   Max.nchar:  17   3rd Qu.:  664.42   
##  Max.   :2023                                     Max.   :18687.50   
##  total_co2_emissions      us_state   
##  Min.   :  0.00      Length   :1388  
##  1st Qu.:  0.30      N.unique :  46  
##  Median :  4.40      N.blank  :   0  
##  Mean   : 17.82      Min.nchar:   4  
##  3rd Qu.: 18.50      Max.nchar:  14  
##  Max.   :538.90
\end{verbatim}

\begin{Shaded}
\begin{Highlighting}[]
\CommentTok{\#numeric summary table}
\NormalTok{blowdown }\SpecialCharTok{\%\textgreater{}\%}
  \FunctionTok{select}\NormalTok{(number\_of\_blowdowns, total\_ch4\_emissions, total\_co2\_emissions) }\SpecialCharTok{\%\textgreater{}\%}
  \FunctionTok{pivot\_longer}\NormalTok{(}\FunctionTok{everything}\NormalTok{(), }\AttributeTok{names\_to =} \StringTok{"variable"}\NormalTok{) }\SpecialCharTok{\%\textgreater{}\%}
  \FunctionTok{group\_by}\NormalTok{(variable) }\SpecialCharTok{\%\textgreater{}\%}
  \FunctionTok{summarise}\NormalTok{(}
    \AttributeTok{n       =} \FunctionTok{sum}\NormalTok{(}\SpecialCharTok{!}\FunctionTok{is.na}\NormalTok{(value)),}
    \AttributeTok{missing =} \FunctionTok{sum}\NormalTok{(}\FunctionTok{is.na}\NormalTok{(value)),}
    \AttributeTok{mean    =} \FunctionTok{mean}\NormalTok{(value, }\AttributeTok{na.rm =} \ConstantTok{TRUE}\NormalTok{),}
    \AttributeTok{median  =} \FunctionTok{median}\NormalTok{(value, }\AttributeTok{na.rm =} \ConstantTok{TRUE}\NormalTok{),}
    \AttributeTok{sd      =} \FunctionTok{sd}\NormalTok{(value, }\AttributeTok{na.rm =} \ConstantTok{TRUE}\NormalTok{),}
    \AttributeTok{min     =} \FunctionTok{min}\NormalTok{(value, }\AttributeTok{na.rm =} \ConstantTok{TRUE}\NormalTok{),}
    \AttributeTok{max     =} \FunctionTok{max}\NormalTok{(value, }\AttributeTok{na.rm =} \ConstantTok{TRUE}\NormalTok{)}
\NormalTok{  )}
\end{Highlighting}
\end{Shaded}

\begin{verbatim}
## # A tibble: 3 x 8
##   variable                n missing  mean median     sd   min    max
##   <chr>               <int>   <int> <dbl>  <dbl>  <dbl> <dbl>  <dbl>
## 1 number_of_blowdowns  1388       0  60.3   12    199.      0  3386 
## 2 total_ch4_emissions  1388       0 621.   164.  1264.      0 18688.
## 3 total_co2_emissions  1388       0  17.8    4.4   39.0     0   539.
\end{verbatim}

The summary above shows the mean, median, and sd of our sample data.

\subsection{Facility Info}\label{facility-info}

\begin{verbatim}
## [1] 56
\end{verbatim}

\begin{verbatim}
## # A tibble: 0 x 2
## # i 2 variables: facility_id <int>, unique_names <int>
\end{verbatim}

\begin{verbatim}
## # A tibble: 0 x 2
## # i 2 variables: facility_name <chr>, unique_ids <int>
\end{verbatim}

A one-to-one match between the Facility Name and Facility ID was
confirmed.

\#Temporal Coverage Looking at years with most/fewest reports. This
should show which years have the most reporting activity. Gaps or sudden
drops in a year could indicate incomplete data or a policy change in EPA
reporting requirements.

\begin{Shaded}
\begin{Highlighting}[]
\FunctionTok{table}\NormalTok{(blowdown}\SpecialCharTok{$}\NormalTok{reporting\_year)}
\end{Highlighting}
\end{Shaded}

\begin{verbatim}
## 
## 2016 2017 2018 2019 2020 2021 2022 2023 
##  140  153  167  184  196  190  183  175
\end{verbatim}

\begin{Shaded}
\begin{Highlighting}[]
\FunctionTok{ggplot}\NormalTok{(blowdown, }\FunctionTok{aes}\NormalTok{(}\AttributeTok{x =}\NormalTok{ reporting\_year)) }\SpecialCharTok{+} \FunctionTok{geom\_bar}\NormalTok{()}
\end{Highlighting}
\end{Shaded}

\pandocbounded{\includegraphics[keepaspectratio]{test-markdown_files/test-markdown_files/figure-latex/unnamed-chunk-4-1.pdf}}

\begin{Shaded}
\begin{Highlighting}[]
\NormalTok{t }\OtherTok{\textless{}{-}} \FunctionTok{as.data.frame}\NormalTok{(}\FunctionTok{table}\NormalTok{(blowdown}\SpecialCharTok{$}\NormalTok{reporting\_year))}
\FunctionTok{colnames}\NormalTok{(t) }\OtherTok{\textless{}{-}} \FunctionTok{c}\NormalTok{(}\StringTok{"Reporting Year"}\NormalTok{, }\StringTok{"Count"}\NormalTok{)}
\FunctionTok{kable}\NormalTok{(t)}
\end{Highlighting}
\end{Shaded}

{\def\LTcaptype{none} % do not increment counter
\begin{longtable}[]{@{}lr@{}}
\toprule\noalign{}
Reporting Year & Count \\
\midrule\noalign{}
\endhead
\bottomrule\noalign{}
\endlastfoot
2016 & 140 \\
2017 & 153 \\
2018 & 167 \\
2019 & 184 \\
2020 & 196 \\
2021 & 190 \\
2022 & 183 \\
2023 & 175 \\
\end{longtable}
}

2020 was year with max reports.

\#Emissions Trends Over Time Shows whether total or average CH4
emissions are going up, down, or staying flat across the industry year
over year.

\begin{Shaded}
\begin{Highlighting}[]
\NormalTok{blowdown }\SpecialCharTok{\%\textgreater{}\%} \FunctionTok{group\_by}\NormalTok{(reporting\_year) }\SpecialCharTok{\%\textgreater{}\%}
  \FunctionTok{summarise}\NormalTok{(}\AttributeTok{mean\_ch4 =} \FunctionTok{mean}\NormalTok{(total\_ch4\_emissions, }\AttributeTok{na.rm=}\NormalTok{T),}
            \AttributeTok{total\_ch4 =} \FunctionTok{sum}\NormalTok{(total\_ch4\_emissions, }\AttributeTok{na.rm=}\NormalTok{T)) }\SpecialCharTok{\%\textgreater{}\%}
  \FunctionTok{ggplot}\NormalTok{(}\FunctionTok{aes}\NormalTok{(}\AttributeTok{x=}\NormalTok{reporting\_year, }\AttributeTok{y=}\NormalTok{total\_ch4)) }\SpecialCharTok{+} \FunctionTok{geom\_line}\NormalTok{()}
\end{Highlighting}
\end{Shaded}

\pandocbounded{\includegraphics[keepaspectratio]{test-markdown_files/test-markdown_files/figure-latex/unnamed-chunk-5-1.pdf}}
hows same thing as above.

\#Top Emitters Will show which facilities are responsible for the most
cumulative CH4 over all years. Useful for understanding whether
emissions are concentrated in a handful of bad actors or spread across
many facilities).

\begin{Shaded}
\begin{Highlighting}[]
\NormalTok{blowdown }\SpecialCharTok{\%\textgreater{}\%} \FunctionTok{group\_by}\NormalTok{(facility\_name) }\SpecialCharTok{\%\textgreater{}\%}
  \FunctionTok{summarise}\NormalTok{(}\AttributeTok{total\_ch4 =} \FunctionTok{sum}\NormalTok{(total\_ch4\_emissions, }\AttributeTok{na.rm=}\NormalTok{T)) }\SpecialCharTok{\%\textgreater{}\%}
  \FunctionTok{slice\_max}\NormalTok{(total\_ch4, }\AttributeTok{n=}\DecValTok{10}\NormalTok{)}
\end{Highlighting}
\end{Shaded}

\begin{verbatim}
## # A tibble: 10 x 2
##    facility_name                                                       total_ch4
##    <chr>                                                                   <dbl>
##  1 Transmission Pipeline Facility, Natural Gas Pipeline Company of Am~   114353.
##  2 Transmission Pipeline - Panhandle Eastern Pipeline Co                  66192.
##  3 Texas Eastern Transmission, LP                                         50772.
##  4 TransCanada / ANR Pipeline                                             45360.
##  5 Transmission Pipeline Facility, El Paso Natural Gas Company, L.L.C.    43229.
##  6 Transmission Pipeline - Florida Gas Transmission                       40881.
##  7 Eastern Gas Transmission and Storage - Transmission Pipelines (TPL)    35714.
##  8 Transmission Pipeline Facility, Tennessee Gas Pipeline Company, L.~    35308.
##  9 Transcontinental Gas Pipe Line Company - Pipeline                      33518.
## 10 TransCanada / Columbia Gas Transmission                                29668.
\end{verbatim}

\#Blowdowns vs.~CH4 Emissions Tests whether more blowdown events at a
facility actually predict higher methane emissions. A strong correlation
would confirm blowdowns are a meaningful driver of leakage; a weak one
suggests other factors (facility size, equipment type) matter more).

\begin{Shaded}
\begin{Highlighting}[]
\FunctionTok{ggplot}\NormalTok{(blowdown, }\FunctionTok{aes}\NormalTok{(}\AttributeTok{x=}\NormalTok{number\_of\_blowdowns, }\AttributeTok{y=}\NormalTok{total\_ch4\_emissions)) }\SpecialCharTok{+}
  \FunctionTok{geom\_point}\NormalTok{(}\AttributeTok{alpha=}\FloatTok{0.4}\NormalTok{) }\SpecialCharTok{+} \FunctionTok{geom\_smooth}\NormalTok{(}\AttributeTok{method=}\StringTok{"lm"}\NormalTok{) }\SpecialCharTok{+}
  \FunctionTok{scale\_y\_log10}\NormalTok{() }\SpecialCharTok{+} \FunctionTok{scale\_x\_log10}\NormalTok{()}
\end{Highlighting}
\end{Shaded}

\begin{verbatim}
## Warning in scale_y_log10(): log-10 transformation introduced infinite values.
\end{verbatim}

\begin{verbatim}
## Warning in scale_x_log10(): log-10 transformation introduced infinite values.
\end{verbatim}

\begin{verbatim}
## Warning in scale_y_log10(): log-10 transformation introduced infinite values.
\end{verbatim}

\begin{verbatim}
## Warning in scale_x_log10(): log-10 transformation introduced infinite values.
\end{verbatim}

\begin{verbatim}
## `geom_smooth()` using formula = 'y ~ x'
\end{verbatim}

\begin{verbatim}
## Warning: Removed 124 rows containing non-finite outside the scale range
## (`stat_smooth()`).
\end{verbatim}

\pandocbounded{\includegraphics[keepaspectratio]{test-markdown_files/test-markdown_files/figure-latex/unnamed-chunk-7-1.pdf}}

\begin{Shaded}
\begin{Highlighting}[]
\FunctionTok{cor}\NormalTok{(blowdown}\SpecialCharTok{$}\NormalTok{number\_of\_blowdowns, blowdown}\SpecialCharTok{$}\NormalTok{total\_ch4\_emissions, }\AttributeTok{use=}\StringTok{"complete.obs"}\NormalTok{)}
\end{Highlighting}
\end{Shaded}

\begin{verbatim}
## [1] 0.2219376
\end{verbatim}

Just becasue you have more blowdown events does not mean you have more
emissions.

\#Geographic Variation by State Reveals which states have the highest
total emissions (useful for identifying regional hotspots; potentiall
states like Texas or Wyoming with heavy pipeline infrastructure will
likely dominate).

\begin{Shaded}
\begin{Highlighting}[]
\NormalTok{blowdown }\SpecialCharTok{\%\textgreater{}\%} \FunctionTok{group\_by}\NormalTok{(us\_state) }\SpecialCharTok{\%\textgreater{}\%}
  \FunctionTok{summarise}\NormalTok{(}\AttributeTok{total\_ch4 =} \FunctionTok{sum}\NormalTok{(total\_ch4\_emissions, }\AttributeTok{na.rm=}\NormalTok{T)) }\SpecialCharTok{\%\textgreater{}\%}
  \FunctionTok{arrange}\NormalTok{(}\FunctionTok{desc}\NormalTok{(total\_ch4))}
\end{Highlighting}
\end{Shaded}

\begin{verbatim}
## # A tibble: 46 x 2
##    us_state     total_ch4
##    <chr>            <dbl>
##  1 Texas          165991.
##  2 Kansas          64303.
##  3 Illinois        48350.
##  4 Iowa            46622.
##  5 Pennsylvania    37053.
##  6 California      36955.
##  7 Louisiana       35935.
##  8 Mississippi     34495.
##  9 Ohio            34368.
## 10 Kentucky        30395.
## # i 36 more rows
\end{verbatim}

Texas is highest emitting state.

\#Panel Data Check of facilities reported across multiple years Do the
same facilities show up across multiple years?.

\begin{Shaded}
\begin{Highlighting}[]
\NormalTok{table\_data }\OtherTok{\textless{}{-}}\NormalTok{ blowdown }\SpecialCharTok{\%\textgreater{}\%}
  \FunctionTok{group\_by}\NormalTok{(facility\_id, facility\_name) }\SpecialCharTok{\%\textgreater{}\%}
  \FunctionTok{summarise}\NormalTok{(}
    \AttributeTok{count          =} \FunctionTok{n}\NormalTok{(),}
    \AttributeTok{years\_reported =} \FunctionTok{paste}\NormalTok{(}\FunctionTok{sort}\NormalTok{(}\FunctionTok{unique}\NormalTok{(reporting\_year)), }\AttributeTok{collapse =} \StringTok{", "}\NormalTok{),}
    \AttributeTok{states         =} \FunctionTok{paste}\NormalTok{(}\FunctionTok{sort}\NormalTok{(}\FunctionTok{unique}\NormalTok{(us\_state)), }\AttributeTok{collapse =} \StringTok{", "}\NormalTok{),}
    \AttributeTok{cumulative\_ch4 =} \FunctionTok{sum}\NormalTok{(total\_ch4\_emissions, }\AttributeTok{na.rm =} \ConstantTok{TRUE}\NormalTok{),}
    \AttributeTok{.groups =} \StringTok{"drop"}
\NormalTok{  ) }\SpecialCharTok{\%\textgreater{}\%}
  \FunctionTok{arrange}\NormalTok{(}\FunctionTok{desc}\NormalTok{(count)) }\SpecialCharTok{\%\textgreater{}\%}
  \FunctionTok{rename}\NormalTok{(}
    \StringTok{"Facility ID"}              \OtherTok{=}\NormalTok{ facility\_id,}
    \StringTok{"Facility Name"}            \OtherTok{=}\NormalTok{ facility\_name,}
    \StringTok{"Count"}                    \OtherTok{=}\NormalTok{ count,}
    \StringTok{"Years Reported"}           \OtherTok{=}\NormalTok{ years\_reported,}
    \StringTok{"States"}                   \OtherTok{=}\NormalTok{ states,}
    \StringTok{"Cumulative CH4 Emissions"} \OtherTok{=}\NormalTok{ cumulative\_ch4}
\NormalTok{  )}
\end{Highlighting}
\end{Shaded}

This table shows one row per facility with the following info: -
Facility ID / Facility Name: identifies the facility - Count: how many
times that facility appears in the dataset (i.e.~how many years it
reported data) - Years Reported: the specific years that facility
submitted emissions data, listed out - States: the state(s) associated
with that facility across its reports - Cumulative CH4 Emissions: the
total methane emissions summed across all years that facility reported.

The table is sorted by Count descending, so facilities that reported the
most years appear at the top to see which facilities have the longest
reporting history and how their total methane output compares to others.

\end{document}
